\section{Background and Motivation}
\label{sec:background}

\subsection{Tiered-memory page placement}

A tiered-memory server presents one virtual address space backed by media with
different performance.  The kernel initially allocates or demotes cold pages
to \slow and promotes frequently accessed pages to \fast.  Linux can infer
temperature through page-table accessed bits, NUMA hint faults, or sampling
facilities such as DAMON.  Research systems add application-transparent
classification~\cite{agarwal2017thermostat}, parallel migration
engines~\cite{yan2019nimble}, and policies tailored to CXL
topologies~\cite{maruf2023tpp}.

All such controllers face a delay.  Let $t_0$ be the first access after a phase
transition, $W$ the sampling window, and $T_m$ the time needed to copy a page.
A reactive policy cannot serve the page from \fast before

\begin{equation}
  t_{\mathrm{ready}} \ge t_0 + W + T_m.
  \label{eq:reactive-delay}
\end{equation}

Reducing $W$ increases sampling overhead and makes classifications noisier;
reducing $T_m$ consumes more copy bandwidth.  Neither option eliminates the
first slow accesses.  \system instead predicts at time $t_p < t_0$ and succeeds
when the \emph{phase exposure} $E=t_0-t_p$ is at least $T_m$.

\subsection{Why tail latency is different}

Capacity services tolerate occasional remote accesses when those accesses
overlap with computation.  Latency-sensitive services often cannot.  A
request may serially traverse an index, wait for an RPC, and then touch a
tenant-specific model.  A few remote misses on this critical chain extend the
request, enlarge queues, and amplify later requests.  This is why optimizing
mean bandwidth or average memory-access time does not necessarily optimize
\ptail latency~\cite{dean2013tail}.

\begin{figure}[t]
  \centering
  \resizebox{0.98\columnwidth}{!}{%
  \begin{tikzpicture}[
    x=0.90cm,y=0.63cm,
    axis/.style={-Latex,thick},
    event/.style={draw,rounded corners=1.5pt,minimum height=0.52cm,
                  align=center,font=\scriptsize},
    timelabel/.style={font=\scriptsize,align=center}
  ]
    \draw[axis] (0,0) -- (7.4,0) node[right,font=\scriptsize] {time};
    \draw[fluxgray,thick] (0.2,1.0) -- (7.0,1.0);
    \node[timelabel,anchor=east] at (0.1,1.0) {reactive};
    \draw[fluxteal,thick] (0.2,2.2) -- (7.0,2.2);
    \node[timelabel,anchor=east] at (0.1,2.2) {\system};

    \node[event,fill=fluxlight] at (1.25,2.2) {phase\\signal};
    \node[event,fill=fluxorange!13] at (3.05,2.2) {copy pages};
    \node[event,fill=fluxteal!15] at (5.65,2.2) {serve from\\\fast};
    \draw[-Latex,fluxteal] (1.75,2.2) -- (2.3,2.2);
    \draw[-Latex,fluxteal] (3.8,2.2) -- (4.85,2.2);

    \node[event,fill=fluxred!12] at (2.70,1.0) {first far\\access};
    \node[event,fill=fluxlight] at (4.25,1.0) {classify};
    \node[event,fill=fluxorange!13] at (5.55,1.0) {copy};
    \node[event,fill=fluxteal!15] at (6.75,1.0) {local};
    \draw[-Latex,fluxgray] (3.3,1.0) -- (3.7,1.0);
    \draw[-Latex,fluxgray] (4.8,1.0) -- (5.0,1.0);
    \draw[-Latex,fluxgray] (6.0,1.0) -- (6.2,1.0);

    \draw[<->,fluxorange,thick] (1.25,2.78) -- (2.70,2.78)
      node[midway,above,font=\scriptsize] {$E$};
    \draw[dashed,black!35] (2.70,0.55) -- (2.70,2.92);
  \end{tikzpicture}%
  }
  \caption{Reactive tiering begins classification after the first far-memory
    access. \system uses phase exposure $E$ to overlap migration with work that
    precedes the phase's memory-intensive region.}
  \Description{Two horizontal timelines compare reactive tiering with
    Flux-Tier. Flux-Tier observes a phase signal and copies pages before the
    first far-memory access; reactive tiering classifies and copies afterward.}
  \label{fig:motivation}
\end{figure}

\subsection{Opportunity and requirements}

We recorded last-level-cache miss samples and instruction pointers for four
services under a 1\,ms observation interval.  Across repeated phases, 71--93\%
of the hottest 2\,MiB page regions reappeared, while a phase-identifying code
sample preceded the first access by 0.4--3.8\,ms.  This interval is long enough
to migrate tens to hundreds of 4\,KiB pages, but not enough to discover them
reactively.  The observations motivate four requirements:

\begin{enumerate}
  \item prediction must be available before the target accesses;
  \item metadata and sampling must be inexpensive at request-scale periods;
  \item inaccurate predictions must not monopolize local capacity or bandwidth;
  \item anonymous services must receive isolation without source changes.
\end{enumerate}
